\section{量子線形方程式問題}
\label{sec:quantum-linear-system-problem}

\subsection{問題の設定}
\label{subsec:qlsp-formulation}

エルミート行列$A\in\mathbb C^{N\times N}$と，
零でないベクトル$\bm{b}\in\mathbb{C}^{N}$に対する線形方程式
\begin{align*}
  A\bm{x}=\bm{b}
\end{align*}
を考える．
行列$A$は逆行列を持つとする．$\bm{b}$の各成分を振幅に記録した量子状態を
\begin{equation}
  \ket{b}
  \coloneq
  \frac{1}{\|\bm{b}\|_2}
  \sum_{k=0}^{N-1}b_k\ket{k}
  \label{eq:qlsp-input-state}
\end{equation}
と定義する．
必要な量子ビット数を$n\coloneq\lceil\log_2N\rceil$とする．
次元$N$が2の累乗でない場合は，使用しない基底状態の振幅を零とし，
行列を
\begin{align*}
  A'
  \coloneq
  A\oplus I_{2^n-N}
\end{align*}
のように広げれば，逆行列を持つという性質を保ったまま$n$量子ビットで表せる．
ここで$A\oplus I_{2^n-N}$は，$A$と恒等行列を対角線上に並べた行列である．

\textbf{量子線形方程式問題}では，古典ベクトル
$\bm{x}=A^{-1}\bm{b}$の全成分を出力しない．
代わりに，解ベクトルの長さを1にそろえた量子状態
\begin{equation}
  \ket{x}
  \coloneq
  \frac{A^{-1}\ket{b}}
       {\|A^{-1}\ket{b}\|_2}
  \label{eq:qlsp-solution-state}
\end{equation}
を生成することを目的とする\cite{shimada2020}．
したがって，HHLアルゴリズムの出力は$N$個の成分を並べた古典データではない．
得られた$\ket{x}$に対して測定を行い，必要な期待値や他の量子状態との内積を評価することが基本となる．

\subsection{固有値分解による解の表現}
\label{subsec:eigendecomposition-of-solution}

行列$A$はエルミートであるため，実数の固有値$\lambda_j$と正規直交固有ベクトル$\ket{u_j}$を用いて
\begin{align}
  A
  &=
  \sum_{j=0}^{N-1}\lambda_j\ket{u_j}\bra{u_j},
  \label{eq:a-spectral-decomposition}\\
  A\ket{u_j}
  &=
  \lambda_j\ket{u_j}
  \label{eq:a-eigenvalue-equation}
\end{align}
と表される．
行列$A$は可逆であるため，すべての固有値は$\lambda_j\neq0$を満たす．

入力状態$\ket{b}$を$A$の固有ベクトルで展開すると
\begin{equation}
  \ket{b}
  =
  \sum_{j=0}^{N-1}\beta_j\ket{u_j},
  \qquad
  \beta_j
  \coloneq
  \braket{u_j|b}
  \label{eq:b-eigenbasis-expansion}
\end{equation}
となる．
このとき，逆行列を作用させた状態は
\begin{equation}
  A^{-1}\ket{b}
  =
  \sum_{j=0}^{N-1}
  \frac{\beta_j}{\lambda_j}\ket{u_j}
  \label{eq:inverse-a-on-b}
\end{equation}
である．
したがって，HHLアルゴリズムの中心となる処理は，
各固有ベクトル成分の振幅$\beta_j$に固有値の逆数$1/\lambda_j$を掛けることである．

\subsection{条件数}
\label{subsec:condition-number}

線形方程式の解が誤差の影響をどの程度受けるかは，行列$A$の条件数に依存する．
スペクトルノルムを用いて，条件数を
\begin{equation}
  \kappa(A)
  \coloneq
  \|A\|_2\|A^{-1}\|_2
  \label{eq:condition-number-definition}
\end{equation}
と定義する．
行列$A$が正定値である場合，最大固有値と最小固有値をそれぞれ
$\lambda_{\max}$，$\lambda_{\min}$とすると
\begin{align*}
  \kappa(A)
  =
  \frac{\lambda_{\max}}{\lambda_{\min}}
\end{align*}
となる．
条件数が大きいほど，小さな固有値に対応する成分は$A^{-1}$によって大きくなる．
その結果，入力や固有値の誤差が解へ大きく影響する．
また，HHLアルゴリズムの成功確率が下がり，より高い精度が必要になる．

線形回帰で$A=X^{\mathsf{T}}X$とした場合を考える．
$X$の列が線形独立ならば
\begin{equation}
  \kappa(X^{\mathsf{T}}X)
  =
  \kappa(X)^2
  \label{eq:normal-equation-condition-number}
\end{equation}
である．正規方程式への書き換えにより，線形回帰へHHLアルゴリズムを使いやすくなる．
一方で，条件数が元の行列の2乗になるため，誤差の影響が大きくなる可能性がある．
